\documentclass[a4paper,11pt]{article}
\usepackage[margin=27mm]{geometry}
\usepackage{fontspec,amsmath,amssymb,amsthm,unicode-math,booktabs}
\setmainfont{STIXGeneral}\setmathfont{STIX Two Math}
\setlength{\emergencystretch}{2em}
\newcommand{\gen}[1]{\langle#1\rangle}
\newtheorem{theorem}{Theorem}\newtheorem{lemma}[theorem]{Lemma}
\title{A counterexample to the six-point case of rank augmentation}
\author{Human Compression project}\date{Private candidate; independent review pending}
\begin{document}\maketitle
\begin{abstract}
We give an explicit rank-three string C-marking of $S_{18}$ for which
splitting one transposition from the middle generator satisfies case (3)
of Theorem 4.1 in the author edition of De Saedeleer--Leemans--Mulpas,
but does not give a rank-four C-marking. Generation of $S_{18}$ follows
from a tree of explicitly conjugate transpositions, and the original
intersection property follows from two dihedral groups. A word in the
last three new marks equals the first new mark, exhibiting the failure.
This does not address cases (1) and (2), the source's small-degree census,
or a complete replacement of that paper.
\end{abstract}

\section{The exact statement being tested}
The source [DLM] is the complete nineteen-page author edition
arXiv:2212.12563v1, 23 December 2022. Its Theorem 4.1 begins on physical
page 5 and continues on page 6. In the notation used here, it starts with
a rank-three string C-marking $(a,q,c)$ of the natural symmetric group
$S_n$, $n\geq5$, with $q$ not a transposition. One writes $q=tb$, where
$t$ is one of its transpositions and neither endpoint of $t$ is moved by
$c$. The proposed new marking is $(t,a,b,c)$. Its subgroup $\gen{a,b}$
is required to have all orbits of size at most three. Case (3) then assumes
that $\gen{t,a,b}$ has an orbit of size six. Unlike cases (1) and (2),
it imposes no imprimitivity or parity condition.

We use right actions and multiply from left to right; thus
$x^g=g^{-1}xg$ sends the support of a transposition $x$ to its image under
$g$. A string C-marking consists of distinct nonidentity involutions,
commuting at nonadjacent positions, with
\[
 \gen{s_i:i\in I}\cap\gen{s_i:i\in J}
 =\gen{s_i:i\in I\cap J}
\]
for every pair of subsets of the index set. We use the standard end
criterion, stated as [DLM, Proposition 5.1] and attributed there to
McMullen--Schulte, Proposition 2E16(a). For a triple, it says that the
intersection of the two rank-two end groups must be the common cyclic
subgroup; the two distinct involutions in each end already form a
rank-two C-marking.

\section{The original marking}
On $\{1,\ldots,18\}$ define
\begin{align*}
 a&=(1\,2)(5\,6)(7\,8)(11\,12)(13\,14)(17\,18),\\
 b&=(2\,3)(4\,5)(7\,9)(10\,11)(14\,15)(16\,17),\\
 c&=(3\,4)(5\,7)(6\,8)(9\,10)(11\,13)(12\,14)(15\,16),\\
 t&=(1\,18),\qquad q=tb.
\end{align*}
The three marks $a,q,c$ are distinct nonidentity involutions, and $a$
commutes with $c$: their common nontrivial supports are the two squares
$\{5,6,7,8\}$ and $\{11,12,13,14\}$. The remaining transpositions of
$a$ and $c$ are disjoint.

\begin{lemma}\label{generation}
The group $\gen{a,q,c}$ is $S_{18}$.
\end{lemma}
\begin{proof}
A direct calculation gives
\[
 aqcq=(1\,5\,8\,11\,15\,12\,7\,6\,2)(4\,9)
       (10\,13\,17\,18\,14).
\]
Consequently $z=(aqcq)^{45}=(4\,9)$. The following table lists seventeen
conjugates of this one transposition. The word in the first column is $g$;
the second column is $z^g$.
\begin{center}
\begin{tabular}{ll@{\qquad}ll}
\toprule
$g$&$z^g$&$g$&$z^g$\\\midrule
$1$&$(4\,9)$&$q$&$(5\,7)$\\
$c$&$(3\,10)$&$qa$&$(6\,8)$\\
$cq$&$(2\,11)$&$cqa$&$(1\,12)$\\
$cqc$&$(2\,13)$&$cqaq$&$(12\,18)$\\
$cqac$&$(1\,14)$&$cqcq$&$(3\,13)$\\
$cqaqa$&$(11\,17)$&$cqacq$&$(15\,18)$\\
$cqcqa$&$(3\,14)$&$cqcqc$&$(4\,11)$\\
$cqaqaq$&$(10\,16)$&$cqcqcq$&$(5\,10)$\\
$cqcqcqa$&$(6\,10)$&&\\\bottomrule
\end{tabular}
\end{center}
The seventeen displayed edges form a connected graph on all eighteen
points. Transpositions on the edges of a connected graph generate its
full symmetric group: along a path, conjugating the next edge
transposition by those already obtained gives a transposition from the
initial vertex to the next vertex. Those transpositions generate all
transpositions. Every displayed transposition belongs to $\gen{a,q,c}$,
so that group contains $S_{18}$ and is therefore equal to it.
\end{proof}

\begin{lemma}\label{oldc}
The original triple $(a,q,c)$ is an irreducible string C-marking,
of type $[6,90]$.
\end{lemma}
\begin{proof}
Put $r=aq$. Direct multiplication gives
\[
 r=(1\,3\,2\,18\,16\,17)(4\,5\,6)(7\,8\,9)
   (10\,11\,12)(13\,15\,14),
\]
so $r$ has order six and
$r^3=(1\,18)(2\,17)(3\,16)$. Also
\[
 bc=(2\,4\,7\,10\,13\,11\,9\,5\,3)(6\,8)
       (12\,14\,16\,17\,15).
\]
The disjoint transposition $t$ commutes with $bc$, so $qc=tbc$ has order
ninety. The two end groups $A=\gen{a,q}$ and $C=\gen{q,c}$ are therefore
dihedral, of orders twelve and one hundred eighty.

The group $C$ preserves each of the following four sets:
\[
 \{1,18\},\quad \{2,3,4,5,7,9,10,11,13\},\quad
 \{6,8\},\quad \{12,14,15,16,17\}.
\]
Every element of $A$ is $r^i$ or $qr^i$. Since $q\in C$, any intersection
larger than $\gen q$ would give a nonidentity rotation $r^i\in C$.
Preservation of $\{1,18\}$ by that rotation forces $i=3$ modulo six.
But $r^3$ sends $2$ to $17$, between two of the displayed $C$-invariant
sets. Thus $A\cap C=\gen q$. The end criterion proves the full
intersection property. The adjacent orders six and ninety also give
irreducibility.
\end{proof}

\section{The source hypotheses and the failure}
The middle involution $q$ is the product of the seven disjoint
transpositions in $t$ and $b$, so it is not itself a transposition.
Replacing its edge $\{1,18\}$ by the new first label yields exactly
$(t,a,b,c)$. Neither $1$ nor $18$ is moved by $c$, so the chosen edge is
not adjacent to any last-label edge. The four new elements are distinct
nonidentity involutions. In addition to $ac=ca$, both $tb=bt$ and
$tc=ct$ hold, so they form a string marking. They generate $S_{18}$
since they generate $a,q,c$. In particular, both the original and the new
group actions are transitive and primitive in the natural degree eighteen.
No fracture graph or additional primitivity hypothesis appears in the
statement being tested.

The six orbits of $B=\gen{a,b}$ are
\[
 \{1,2,3\},\ \{4,5,6\},\ \{7,8,9\},\
 \{10,11,12\},\ \{13,14,15\},\ \{16,17,18\}.
\]
Each is the connected path of one $a$-edge and one $b$-edge.
The subgroup $H=\gen{t,a,b}$ joins the first and last of these triples
and fixes their union setwise. That union is connected under $t,a,b$,
so it is one orbit of size six. The remaining four orbits still have
size three. Thus both orbit hypotheses of [DLM, Theorem 4.1(3)] hold
literally. Even irreducibility of the proposed new marking holds, since
its adjacent orders are $4,3,90$.

Nevertheless its intersection property fails. The displayed cycle form
of $bc$ gives $(bc)^{45}=(6\,8)$. Set
\[
 w=a(bc)^2ba(cb)^8a.
\]
Multiplying the given permutations shows $6^w=1$ and $8^w=18$.
Consequently
\begin{equation}\label{defect}
 ((bc)^{45})^w=(1\,18)=t\in\gen{a,b,c}.
\end{equation}
The index sets of the first mark $t$ and the last three marks are disjoint,
but their generated subgroups have the nonidentity element $t$ in common.
This is an explicit violation of the defining intersection property.

\begin{theorem}
Theorem 4.1(3) of [DLM], in the exact author edition specified above,
is false as stated. The displayed eighteen-point example satisfies all
of its hypotheses and fails its conclusion.
\end{theorem}
\begin{proof}
Lemmas~\ref{generation} and~\ref{oldc} verify the entire input marking.
The preceding orbit and edge checks verify the additional hypotheses.
Equation~\eqref{defect} disproves the conclusion.
\end{proof}

The last paragraph of the source proof on physical page 16 uses the cube
structure of the six-point action to conclude that the two end groups
intersect only in $B$. That local dependence is insufficient here.
Indeed, \eqref{defect} implies $q=tb\in\gen{a,b,c}$; together with
Lemma~\ref{generation} it follows that $\gen{a,b,c}=S_{18}$.
Thus $H\cap\gen{a,b,c}=H$, which is strictly larger than $B$.
No disproof of cases (1) or (2), no complete source recovery, and no
new conclusion about the distinct high-rank stability theorem is claimed.

\section{Verification and attribution}
All ingredients above reduce to finite permutation identities, the
connected-graph transposition argument, and the explicitly retained end
criterion. The complete program \texttt{check\_case.py} verifies these
identities directly on eighteen-entry permutation tuples. It enumerates
only the proper subgroups of the old rank-three marking, of order at most
180, and checks all sixty-four old parabolic pairs. It never enumerates
$S_{18}$: the seventeen-edge certificate proves that group identification.
It also checks the six middle orbits, the six-point end orbit, the splitting
conditions, and the explicit word \eqref{defect}.

The source theorem and intended local cube argument are due to
De Saedeleer, Leemans and Mulpas. The elementary end criterion is retained
from McMullen--Schulte as cited by that source. This candidate has no
end-to-end Lean certificate; the complete mathematical proof above and
its transparent finite replay await distinct review. It is not a numerical
compression claim.

\begin{thebibliography}{9}
\bibitem{DLM} J. De Saedeleer, D. Leemans and J. Mulpas,
\emph{A rank augmentation theorem for rank 3 string C-group representations
of the symmetric groups}, arXiv:2212.12563v1, 23 December 2022,
Theorem 4.1 and \S5.3, physical pages 5--6 and 14--16.
\end{thebibliography}
\end{document}
